\clearpage\fancyhead[L]{\small PALLAS.}\hypertarget{native-384-lower}{}\label{native-384-lower}
\begin{center}\textbf{PALLAS.}\par\vspace{4pt}\hrule\vspace{4pt}\fontsize{10}{12}\selectfont Göttingische gelehrte Anzeigen. Stück 199. Seite 1985\;..\;1990. 12. Dezember 1814.\par\vspace{4pt}\hrule\vspace{4pt}\end{center}

\par\indent Die in unseren Blättern schon öfter erwähnte Arbeit von Professor \textsc{Gauss} über die verwickelte Theorie der Bewegungen der Pallas hat durch die Beobachtungen der Opposition dieses Planeten im Oktober dieses Jahres, der \emph{zehnten} seit seiner Entdeckung, erneut eine schöne Bestätigung erhalten. Wir teilen hier diese Beobachtungen und die \emph{zunächst} daraus abgeleiteten Ergebnisse mit und behalten die weiteren Folgerungen für die Theorie einer anderen Gelegenheit vor. Herr \textsc{Nicolai}, Adjunkt an der Seeberger Sternwarte, hatte bei seiner vor mehr als einem Jahr angestellten Vorausberechnung der Bewegung für die diesmalige Sichtbarkeit nur die oskulierenden Elemente der neunten Opposition zugrunde gelegt (was zum Auffinden des Planeten natürlich vollkommen ausreichen musste). Professor \textsc{Gauss} hielt es daher für interessant, noch vor der Wiederkehr des Planeten zu bestimmen, wie diese Elemente durch die fortwährende Einwirkung des Jupiters von der neunten bis zur zehnten Opposition zu berichtigen waren. Er ließ diese nicht unbedeutende Arbeit schon zu Beginn dieses Jahres unter seiner Leitung von Herrn \textsc{Möbius} ausführen, der sich damals bei uns mit ausgezeichnetem Eifer mathematischen und astronomischen Studien widmete. Die Ergebnisse dieser Berechnung sind ausführlicher bereits im soeben erschienenen astronomischen Jahrbuch für 1817 abgedruckt. Zum Vergleich mit den tatsächlichen Beobachtungen geben wir hier nur die Vorausbestimmung der Opposition selbst wieder.

\begin{center}25. Oktober 1814. $12^{\mathrm u}\;39^{\mathrm m}\;50^{\mathrm s}5$ Mittlere Zeit in Göttingen\end{center}
\begin{center}\fontsize{10}{12}\selectfont\setlength{\tabcolsep}{3pt}\renewcommand{\arraystretch}{1.15}\begin{tabular}{lcr}
Wahre Länge & $.\;.\;.\;.\;.\;.\;.\;.\;.\;.\;.\;.\;.$ & $31^\circ\;58'\;28''4$\\
Heliozentrische Breite & $.\;.\;.\;.\;.\;.\;.\;.\;.\;.$ & $23\;38\;32.7\;\mathrm{S.}$\\
\end{tabular}\end{center}

\SourceLeaf{385}\Needspace{10\baselineskip}\par\indent An der hiesigen Sternwarte wurde mit dem Kreismikrometer die folgende Beobachtung gemacht.

\begin{center}\fontsize{9}{11}\selectfont\setlength{\tabcolsep}{3pt}\renewcommand{\arraystretch}{1.15}\begin{tabular}{c|ccc|ccc|ccc}
1814 & \multicolumn{3}{c|}{\shortstack{Mittlere Zeit\\in Göttingen}} & \multicolumn{3}{c|}{\shortstack{Scheinbare\\Rektaszension}} & \multicolumn{3}{c}{\shortstack{Scheinbare\\Deklination}} \\
\hline
Sept.\ 16 & $12^{\mathrm u}$ & $1^{\mathrm m}$ & $1^{\mathrm s}$ & $46^{\circ}$ & $16'$ & $23''4$ & $11^{\circ}$ & $4'$ & $20''9$ S. \\
\end{tabular}\end{center}

\par\indent Dadurch bestätigte sich die Richtigkeit der Ephemeride. Da kein Zweifel bestand, dass der Planet an anderen Sternwarten mit fest aufgestellten Instrumenten beobachtet werden konnte, wäre es zwecklos gewesen, die Beobachtungen mit dem Kreismikrometer fortzusetzen. Die obige Beobachtung stützte sich übrigens auf zahlreiche Vergleiche und war im Tagebuch als sehr gut bezeichnet.

\Needspace{19\baselineskip}\par\indent An der \emph{Seeberger} Sternwarte beschränkte man sich mit Recht auf Beobachtungen mit dem Meridianfernrohr, da der Planet zu lichtschwach war, um vom \textsc{Dollond}schen Quadranten zuverlässige Deklinationen erwarten zu können. Herr von \textsc{Lindenau} und Herr \textsc{Nicolai} beobachteten die folgenden Rektaszensionen:

\begin{center}\fontsize{9}{11}\selectfont\setlength{\tabcolsep}{3pt}\renewcommand{\arraystretch}{1.15}\begin{tabular}{c|ccc|ccc}
1814 & \multicolumn{3}{c|}{\shortstack{Mittlere Zeit\\auf Seberg}} & \multicolumn{3}{c}{\shortstack{Scheinbare\\Rektaszension}} \\
\hline
Sept.\ 17  & $15^{\mathrm u}$ & $19^{\mathrm m}$ & $41^{\mathrm s}5$ & $46^{\circ}$ & $18'$ & $26''2$ \\
18         & 15 & 15 & 52.0 & 46 & 20 & 2.8 \\
19         & 15 & 12 & 1.0  & 46 & 21 & 17.0 \\
20         & 15 & 8  & 8.0  & 46 & 22 & 0.7 \\
Okt.\ 12 & 13 & 36 & 19.2 & 45 & 2  & 4.5 \\
13         & 13 & 31 & 52.0 & 44 & 54 & 14.6 \\
15         & 13 & 22 & 53.1 & 44 & 37 & 25.8 \\
21         & 12 & 55 & 28.0 & 43 & 39 & 51.6 \\
28         & 12 & 22 & 45.3 & 42 & 21 & 48.5 \\
\end{tabular}\end{center}

\Needspace{14\baselineskip}\par\indent Die Beobachtungen, die Herr \textsc{Schumacher} an der \emph{Mannheimer} Sternwarte mit dem acht Fuß messenden Mauerquadranten angestellt hatte, wurden vollständig im Original eingesandt und hier von Herrn \textsc{Encke} mit größter Sorgfalt reduziert. Daraus ergab sich Folgendes:

\begin{center}\fontsize{9}{11}\selectfont\setlength{\tabcolsep}{3pt}\renewcommand{\arraystretch}{1.15}\begin{tabular}{c|ccc|ccc|ccc}
1814 & \multicolumn{3}{c|}{\shortstack{Mittlere Zeit\\in Mannheim}} & \multicolumn{3}{c|}{\shortstack{Scheinbare\\Rektaszension}} & \multicolumn{3}{c}{\shortstack{Scheinbare\\Deklination}} \\
\hline
Okt.\ 12 & $13^{\mathrm u}$ & $36^{\mathrm m}$ & $17^{\mathrm s}4$ & $45^{\circ}$ & $2'$  & $0''3$  & $19^{\circ}$ & $12'$ & $38''4$ S. \\
15         & 13 & 22 & 51.5 & 44 & 37 & 23.7 & 20 & 6  & 54.6 \\
21         & 12 & 55 & 27.0 & 43 & 39 & 57.5 & 21 & 49 & 52.1 \\
22         & 12 & 50 & 48.3 & 43 & 29 & 15.3 & 22 & 6  & 3.8  \\
\end{tabular}\end{center}

\par\indent Bei den beiden letzten Beobachtungen war der Planet so lichtschwach, dass er sich nur mit Mühe erkennen ließ.

\Needspace{18\baselineskip}\par\indent An der neuen \emph{Königsberger} Sternwarte (von deren großer Tätigkeit unsere Blätter schon mehrere Beispiele gegeben haben) machte Professor \textsc{Bessel} die folgenden Beobachtungen:

\begin{center}\fontsize{9}{11}\selectfont\setlength{\tabcolsep}{3pt}\renewcommand{\arraystretch}{1.15}\begin{tabular}{c|ccc|ccc|ccc}
1814 & \multicolumn{3}{c|}{\shortstack{Mittlere Zeit\\in Königsberg}} & \multicolumn{3}{c|}{\shortstack{Scheinbare\\Rektaszension}} & \multicolumn{3}{c}{\shortstack{Scheinbare\\Deklination}} \\
\hline
Okt.\ 14 & $13^{\mathrm u}$ & $27^{\mathrm m}$ & $30^{\mathrm s}6$ & $44^{\circ}$ & $46'$ & $15''3$ & $19^{\circ}$ & $48'$ & $30''8$ S. \\
15         & 13 & 23 & 0.4  & 44 & 37 & 38.6 & 20 & 6  & 20.0 \\
18         & 13 & 9  & 22.9 & 44 & 10 & 8.3  & 20 & 58 & 58.9 \\
19         & 13 & 4  & 48.3 & 44 & 0  & 25.9 & 21 & 15 & 52.6 \\
Nov.\ 1    & 12 & 3  & 57.3 & 41 & 33 & 59.9 & 24 & 31 & 29.8 \\
2          & 11 & 59 & 12.6 & 41 & 21 & 46.2 & 24 & 44 & 9.7  \\
3          & 11 & 54 & 27.4 & 41 & 9  & 24.1 & 24 & 56 & $10\;6$ \\
4          & 11 & 49 & 42.0 & 40 & 56 & 59.6 & 25 & 7  & 41.8 \\
\end{tabular}\end{center}

\SourceLeaf{386}\par\indent Auch dieser ausgezeichnete Beobachter beklagt sich über die Schwierigkeiten, die sich bei den Beobachtungen aus der Lichtschwäche des Planeten ergaben.

\Needspace{31\baselineskip}\par\indent Die folgende Übersicht zeigt die Ergebnisse des Vergleichs all dieser Beobachtungen mit den zuletzt gefundenen Elementen. (Dabei ist nur zu bemerken, dass die verwendete Epoche der mittleren Länge um einige Sekunden abweicht. Das liegt an einem Verfahren von Professor \textsc{Gauss}, das ihm eigen ist und an anderer Stelle ausführlicher begründet werden soll.)

\begin{center}
\fontsize{10}{12}\selectfont\setlength{\tabcolsep}{5pt}\renewcommand{\arraystretch}{1.15}\begin{tabular}{l|cc|l}
 & \multicolumn{2}{c|}{Unterschied} & \\
1814 & Rektasz. & Dekl. & Beobachter \\
\hline
Sept.\ 16 & $-9''2$ & $+5''9$ & Gauss \\
17 & $-1.6$ & & von Lindenau \\
18 & $-6.2$ & & von Lindenau \\
19 & $-10.7$ & & von Lindenau \\
20 & $-7.5$ & & von Lindenau \\
Okt.\ 12 & $-4.9$ & & von Lindenau \\
12 & $-3.6$ & $+2.2$ & Schumacher \\
13 & $-8.8$ & & von Lindenau \\
14 & $-9.4$ & $+12.0$ & Bessel \\
15 & $-5.2$ & $+6.7$ & Bessel \\
15 & $-6.5$ & & von Lindenau \\
15 & $-7.7$ & $+5.7$ & Schumacher \\
18 & $-2.7$ & $+13.2$ & Bessel \\
19 & $-4.0$ & $+2.4$ & Bessel \\
21 & $-2.2$ & & von Lindenau \\
21 & $-12.0$ & $+4.5$ & Schumacher \\
22 & $-0.4$ & $-1.3$ & Schumacher \\
28 & $+1.8$ & & von Lindenau \\
Nov.\ 1 & $-0.4$ & $+8.5$ & Bessel \\
2 & $-5.1$ & $+10.6$ & Bessel \\
3 & $-5.9$ & $-2.8$ & Bessel \\
4 & $-7.9$ & $-21.9$ & Bessel \\
\end{tabular}
\end{center}

\Needspace{10\baselineskip}\par\indent Für die \emph{zehnte Opposition} selbst ergibt sich schließlich aus diesen Beobachtungen Folgendes:

\begin{center}25. Oktober 1814. $12^{\mathrm u}\;33^{\mathrm m}\;22^{\mathrm s}$ Mittlere Zeit in Göttingen\end{center}
\begin{center}\fontsize{10}{12}\selectfont\setlength{\tabcolsep}{3pt}\renewcommand{\arraystretch}{1.15}\begin{tabular}{lcr}
Wahre Länge & $.\;.\;.\;.\;.\;.\;.\;.\;.\;.\;.\;.$ & $31^\circ\;58'\;11''3$\\
Geozentrische Breite & $.\;.\;.\;.\;.\;.\;.\;.\;.\;.$ & $37\;20\;53.2\;\mathrm{Südl.}$\\
\end{tabular}\end{center}

\Needspace{21\baselineskip}\par\indent Der genaue Vergleich dieses Ergebnisses mit den zuletzt gefundenen und um die Jupiterstörungen gehörig berichtigten Elementen ergab so geringe Unterschiede, dass die Grundelemente danach keiner weiteren Berichtigung bedürfen. Wir geben daher nur noch die folgende Übersicht darüber, wie gut alle bisher beobachteten Oppositionen mit der Theorie von Herrn \textsc{Gauss} übereinstimmen.

\begin{center}
\fontsize{10}{12}\selectfont\setlength{\tabcolsep}{5pt}\renewcommand{\arraystretch}{1.15}\begin{tabular}{ll|cc}
 & & \multicolumn{2}{c}{Unterschied} \\
Oppos. & & Mittl.\ Länge & Helioz. Breite \\
\hline
I & 1803 & $-1''3$ & $-5''3$ \\
II & 1804 & $-12.9$ & $-11.9$ \\
III & 1805 & $-20.4$ & $-0.4$ \\
IV & 1807 & $+2.5$ & $-6.8$ \\
V & 1808 & $+23.5$ & $+9.0$ \\
VI & 1809 & $+21.5$ & $-3.1$ \\
VII & 1811 & $+14.1$ & $-7.2$ \\
VIII & 1812 & $-16.1$ & $+4.0$ \\
IX & 1813 & $-10.9$ & $-0.0$ \\
X & 1814 & $-14.0$ & $+2.4$ \\
\end{tabular}
\end{center}
